% Generated by build/make_tables.py from book_fixed_notionals.csv. Do not edit.

\begin{table}[tbp]
\centering
\small
\caption{The fixed proxy book after calibration to the FY2019 disclosed risk split}
\label{tab:book}
\begin{tabular}{@{}lllrr@{}}
\toprule
Factor & Category & Instrument & Maturity (y) & Notional (\$mm) \\
\midrule
us\_2y & Interest rates & Par coupon bond & 2 & 8,403 \\
us\_5y & Interest rates & Par coupon bond & 5 & 2,467 \\
us\_10y & Interest rates & Par coupon bond & 10 & 1,302 \\
us\_30y & Interest rates & Par coupon bond & 30 & 553 \\
eur\_2y & Interest rates & Zero-coupon bond & 2 & 18,162 \\
eur\_5y & Interest rates & Zero-coupon bond & 5 & 3,941 \\
eur\_10y & Interest rates & Zero-coupon bond & 10 & 1,427 \\
\addlinespace
equity\_us & Equity prices & Linear in factor & --- & 604 \\
equity\_europe & Equity prices & Linear in factor & --- & 478 \\
equity\_japan & Equity prices & Linear in factor & --- & 444 \\
equity\_em & Equity prices & Linear in factor & --- & 442 \\
\addlinespace
fx\_eurusd & Currency rates & Linear in factor & --- & 788 \\
fx\_usdjpy & Currency rates & Linear in factor & --- & 734 \\
fx\_gbpusd & Currency rates & Linear in factor & --- & 701 \\
\addlinespace
commodity\_gold & Commodity prices & Linear in factor & --- & 443 \\
commodity\_oil & Commodity prices & Linear in factor & --- & 165 \\
commodity\_base\_metals & Commodity prices & Linear in factor & --- & 341 \\
\bottomrule
\end{tabular}
\notetext{Notionals are the output of the sizing described in Section~\ref{sec:method}, not an estimate of anything Goldman holds. Sleeve totals are \$36,255mm of rates, \$1,968mm of equity, \$2,223mm of currency and \$949mm of commodity exposure, \$41,395mm in all. The rate sleeve dominates on notional and not on risk: a two-year note carries a fraction of the daily risk of a dollar of crude oil fund, which is the whole reason sizing targets risk rather than money. The dollar level is a normalisation and carries no information; only the ratios between sleeves do.}
\end{table}
